%=======================================================================
%  ICCJ 2026 - 3-page conference submission
%  Template: IEEEtran (conference option). Compile with pdfLaTeX.
%=======================================================================
\documentclass[conference]{IEEEtran}
\IEEEoverridecommandlockouts

\usepackage{cite}
\usepackage{amsmath,amssymb}
\usepackage{graphicx}
\usepackage{booktabs}
\usepackage{url}
\usepackage[T1]{fontenc}

\begin{document}

\title{Provenance-Locked Measurement of Flood Extent and\\
District News Visibility: Pipeline Design and\\
Feasibility Findings for 190 Indian Flood Events}

\author{\IEEEauthorblockN{Seongmin Roh}
\IEEEauthorblockA{\textit{Affiliation (to be completed)}\\
Seoul, Republic of Korea\\
smroh0509@gmail.com}}

\maketitle

\begin{abstract}
Whether local news visibility of a disaster depends on how urban the affected
area is can only be tested if physical severity is held fixed at the same
spatial unit as the news measurement. We present CVND, a reproducible pipeline
that joins satellite-observed flood extent, Census~2011 urbanization and
GDELT-indexed news counts at the flood-event~$\times$~district level for India
(190 EM-DAT parent events, 2015--2025, 1{,}630 event--district rows, 556
districts). Its design principle is that measurement provenance outranks sample
size: every satellite area carries the hash of a frozen measurement
specification, changing that specification invalidates caches by construction,
and a pre-specified estimability gate reports \emph{not estimable} rather than
fitting an underpowered model. We report three audit findings obtained under
this discipline. First, 458 of 464 unmatched district areas-of-interest (98.7\%)
were recovered, but only by mixing boundary vintages, which makes district areas
non-additive: 422 same-event polygon overlaps across 101 events. Second,
Sentinel-2 L2A is unavailable for 63.5\% of probed district-events and the
failures are concentrated before 2019; an L1C + \textit{s2cloudless} path
recovers 36.4\% of them, at lower median usable tile fraction. Third, a
507{,}076-article state-level GDELT corpus yields district-explicit evidence for
only 762 rows (0.6\% of in-window candidates), so state-level coverage cannot be
down-scaled to districts. Under the pre-specified gates the primary
negative-binomial models are currently not estimable ($N=0$); we report that
status rather than relax the specification.
\end{abstract}

\begin{IEEEkeywords}
disaster informatics, flood mapping, media coverage, GDELT, Sentinel-1/2,
reproducible pipelines, measurement provenance
\end{IEEEkeywords}

%=======================================================================
\section{Introduction}
%=======================================================================
Disaster relief and political attention track news coverage rather than physical
severity alone~\cite{eisensee2007}. If coverage is systematically thinner where
fewer people live in towns, then rural populations are disadvantaged twice: once
by the flood and again by the response it attracts. Satellite archives now make
the physical side of that comparison measurable at scale~\cite{tellman2021}, but
testing the claim requires comparing districts \emph{at similar observed flood
extent}---which means the physical measurement and the news measurement must
exist at the same administrative unit.

Most existing comparisons operate at country or state level, where a single
coverage count is attributed to a heterogeneous area. India's states routinely
contain both metropolitan and predominantly rural districts, so a state-level
comparison cannot separate the two. We therefore target the district (admin-2)
unit throughout, and treat any state-level quantity as inadmissible evidence for
a district observation.

This paper reports the construction, not the conclusion. We contribute a frozen,
hash-identified measurement specification that makes every satellite area
self-invalidating when its definition changes (Sec.~\ref{sec:spec}); a
conservative boundary recovery procedure carrying vintage provenance and an
explicit non-additivity constraint (Sec.~\ref{sec:boundary}); a Sentinel-2
availability probe establishing when an optical-primary measurement is feasible
(Sec.~\ref{sec:s2}); a negative result on corpus reuse (Sec.~\ref{sec:news});
and a pre-specified estimability gate, demonstrated by reporting a currently
unestimable primary model (Sec.~\ref{sec:status}).

%=======================================================================
\section{Study Design}
%=======================================================================
\subsection{Observation unit and event registry}
The primary key is the flood event~$\times$~district. \texttt{event\_id} is the
parent EM-DAT~\cite{emdat} event identifier, \texttt{source\_record\_id} is the
official DisNo., and \texttt{event\_district\_id} is a deterministic district
key. The registry holds 190 parent events drawn from 72 EM-DAT source records,
expanded to 1{,}630 event--district rows over 556 districts and 31 states, with
onsets from 2015-03-20 to 2025-10-03. Structured GADM/Admin-Unit fields are
preferred; conservatively parsed free-text location evidence stays separately
identifiable, and unresolved districts remain audit rows rather than being
replaced by a state capital or a fuzzy name match. Onset precision is retained:
244 rows (15.0\%) have month-imputed onsets, so their 14-day news windows carry
timing uncertainty.

\subsection{Frozen measurement specification}\label{sec:spec}
Every satellite area is produced under one \texttt{MeasurementSpec}: a post
window $[t_0,\,t_0+14\,\mathrm{d})$ matching the news window, a pre-window
median over $[t_0-30\,\mathrm{d},\,t_0)$, a maximum-water post composite, one
eligibility mask applied on every sensor path (JRC permanent
water~\cite{pekel2016}, slope $<5^\circ$, India LSIB), pixel-area sums at
10\,m~\cite{gorelick2017}, and an identical new-water definition for
Sentinel-1 Otsu thresholding~\cite{torres2012,otsu1979}, Sentinel-2
NDWI~\cite{drusch2012,mcfeeters1996} and the satellite image time-series (SITS)
path.

The specification is hashed (\texttt{spec\_version} \texttt{fs1-3016a20fd3ad})
and that hash is stamped on every Track-A row, HDF5 patch, score archive, merged
row and area-table row. Editing the specification therefore invalidates all
satellite caches by construction: stale rows are discarded, never silently
reused. The same mechanism applies to the event registry, whose fingerprint is
carried by every derived artifact and every news collection manifest.

\subsection{Pre-specified inference and estimability gates}
The primary model is a negative binomial (NB2) count
regression~\cite{hilbe2011}:
\begin{align}
Y_{ed} &\sim \mathrm{NB2}(\mu_{ed}, \alpha),\nonumber\\
\log \mu_{ed} &= \beta_0 + \beta_{f}\log\!\left(1 + A_{ed}\right)
                        + \beta_{u} u_{d},
\end{align}
where $A_{ed}$ is flood area in km$^2$, $u_d$ is the Census~2011 district urban
population share~\cite{census2011}, and $\mathrm{Var}(Y_{ed}) = \mu_{ed} +
\alpha\mu_{ed}^2$ with $\alpha$ estimated by maximum likelihood rather than
fixed. Model~1 omits $u_d$; Model~3 adds year fixed effects; all three use the
same complete-case sample. The urbanization effect is reported as
$\exp(0.1\,\beta_u)$ per 10 percentage points, so no arbitrary urban/rural
cutoff enters the specification.

Crucially, the numerical safeguards are fixed in advance and do not depend on
the direction of the result. A model is estimated only with $N \ge 20$, at least
five observations per estimated parameter, a full-rank design and stable
convergence; otherwise the pipeline emits a typed \emph{not estimable} status.
State-clustered covariance with a finite-sample correction and $t(G-1)$
inference~\cite{cameron2015} is added only when $G \ge 20$ and $N > 2G$;
otherwise the independence assumption is stated as an explicit limitation.
Synthetic exposure proxies are excluded by design: area~$\times$~average-density
is never used as an exposed-population surrogate.

%=======================================================================
\section{Pipeline}
%=======================================================================
The runner executes eleven stages: event registry construction; Census
covariates; district AOI resolution; Sentinel-1/2 measurement (Track~A) with
optional SITS patch preparation (Track~B); local SITS inference; candidate
routing; flood-area table assembly; district news collection; article
classification; identity-checked joining; and coverage analysis. Each stage
refuses to manufacture observations: a successful query with zero matching
candidates yields an observed zero, whereas an unexecuted, failed or incomplete
stage yields missing. An offline \texttt{--dry-run} mode prints the command plan
and audits input and cache schemas without external calls, and the repository
carries 218 unit tests.

%=======================================================================
\section{Feasibility and Audit Findings}
%=======================================================================
\subsection{District boundary recovery and non-additivity}\label{sec:boundary}
Resolving a unique polygon per district is the binding constraint on the
physical measurement. An initial FAO GAUL~2015 level-2 lookup left 464 rows
unmatched. A manifest-driven recovery reconnected 458 of them (98.7\%) using
pinned geoBoundaries gbOpen India ADM2 polygons~\cite{runfola2020} for 173
district names and five reviewed OpenStreetMap relations, raising the cache to
1{,}549 \textsc{ok} / 75 \textsc{no\_imagery} / 6 \textsc{error}. Matching is
explicit: state-scoped aliases and manually reviewed renames only, with no fuzzy
nearest-name fallback and no substitution of a parent district or state.
Candidate polygons were checked for validity, unique identity, $\ge 90\%$
overlap with their named state layer and no $>1\%$ overlap with a neighbour in
the same layer. Six identities remain deliberately unresolved rather than
approximated, including a supplemental polygon (Shamator,
$\approx$193\,km$^2$ in OSM versus 469\,km$^2$ officially) rejected as
incomplete.

The recovery buys coverage at a documented price. The resulting dataset mixes
reference vintages---GAUL~2015 parents alongside a 2021 ADM2 layer---so older
parent polygons can overlap newer child districts. An audit found 422
same-event overlapping pairs across 101 events, 419 of them cross-source.
District flood areas therefore \emph{must not be summed} into event totals; a
valid event total requires a common non-overlapping geography or a union of
pixel-level flood masks. Each row therefore carries \texttt{aoi\_source},
\texttt{geometry\_id}, \texttt{aoi\_match\_method},
\texttt{aoi\_boundary\_year} and an explicit
\texttt{aoi\_historical\_boundary\_verified} flag, because a successful match
means a polygon was deliberately selected---not that it equals the boundary in
force on the event date.

\subsection{Sentinel-2 availability for district-scale SITS}\label{sec:s2}
Before committing to an optical-primary route we probed 495 registry rows, of
which 255 were AOI-matched, under the frozen specification
(Table~\ref{tab:s2}). Surface-reflectance (L2A) imagery is unavailable for
63.5\% of probed district-events, dominated by absent post-window imagery
(148 of 162 failures). The failures are not random in time: 95.1\% of them fall
in years up to 2018, consistent with the L2A archive's start, and the yearly
unavailability rate falls from 1.00 (2015--2016) to 0.00 for 2021--2022 and
2024--2025. An L1C path with \textit{s2cloudless} probability $<40$ lowers
unavailability to 40.4\% and rescues 36.4\% of the L2A failures, but its median
usable tile fraction is 0.095 against 0.283 for L2A---more scenes, thinner
scenes. The pre-specified gate accordingly returns \texttt{consider\_l1c}
rather than an unconditional switch, and records that L1C still exceeds the
0.30 maximum unavailability target. The probe also sizes the commitment: a
planned-layout L1C build is 119.7\,GB of HDF5 over 584{,}536 retained tiles and
roughly 59 GPU-hours, against 92.3\,GB and 40 hours for L2A.

\begin{table}[t]
\caption{Sentinel-2 availability probe (255 AOI-matched district-events)}
\label{tab:s2}
\centering
\begin{tabular}{lrr}
\toprule
 & L2A (SR) & L1C + \textit{s2cloudless} \\
\midrule
Unavailable rows          & 162 (63.5\%) & 103 (40.4\%) \\
\quad no post imagery     & 148 & 91 \\
\quad no pre imagery      & 10  & 5  \\
\quad no retained tiles   & 4   & 5  \\
\quad no clear baseline   & 0   & 2  \\
Median usable tile frac.  & 0.283 & 0.095 \\
Retained tiles            & 450{,}358 & 584{,}536 \\
Planned HDF5 volume       & 92.3\,GB & 119.7\,GB \\
Est. inference hours      & 40.3 & 59.4 \\
\bottomrule
\end{tabular}
\end{table}

\subsection{State news corpora do not transfer to districts}\label{sec:news}
News counts use historical GDELT GKG metadata~\cite{leetaru2013} within a fixed
window $t_0 \le t_{\text{article}} < t_0 + 14$\,d (UTC). Explicit district
location evidence with state disambiguation is required; a state mention alone
cannot create district coverage. Within a district, one URL is assigned to its
nearest eligible onset, and the same URL may count for several districts it
explicitly names. We tested whether an existing state-level corpus could
substitute for district collection. Of 507{,}076 archived state-event article rows, 127{,}252 (25.1\%)
passed event identity, state and window selection, but only 762 could be
provisionally assigned to a district from available evidence---0.15\% of the
corpus and 0.6\% of in-window candidates. The remaining 126{,}557 in-window rows
carry no resolvable district. State-level corpora are therefore not a shortcut:
district-explicit collection or per-article adjudication is unavoidable, and
counts obtained this way remain candidates, not research observations.

\subsection{Current estimability status}\label{sec:status}
Table~\ref{tab:qc} gives stage completion against analysis admission. District
identity extraction succeeds for all 1{,}630 rows and Census matching for 1{,}488
(91.3\%). The satellite cache reached 1{,}549 resolved rows, yet zero rows are
admitted to analysis: the event registry was regenerated, its fingerprint
changed, and the invalidation mechanism of Sec.~\ref{sec:spec} correctly
discarded the derived satellite artifacts that had been produced against the
previous registry. District news collection is likewise recorded as incomplete
rather than zero.

Consequently $N=0$ and all pre-specified models return typed failures
(\texttt{insufficient sample: N=0 < 20} for Models~1--2, fewer than two years
for Model~3, fewer than two satellite sources for the source fixed-effect
sensitivity). No estimate of $\beta_f$ or $\beta_u$ is reported, and the
conclusion field reads \emph{no empirical conclusion is available}. Selection
QC confirms the exclusion is uniform---411/411 high, 399/399 medium, 692/692 low
and 128/128 unknown-urbanization rows---so the present attrition is a pipeline
state, not a differential one. Among rows that do carry covariates, district
urban population share has mean 0.221 and median 0.150 ($n=1{,}502$),
confirming that the registry spans a wide urbanization range once measurement
resumes.

\begin{table}[t]
\caption{Stage completion vs.\ analysis admission (1{,}630 rows)}
\label{tab:qc}
\centering
\begin{tabular}{lrrr}
\toprule
Stage & Resolved & Rate & Admitted \\
\midrule
District identity extraction   & 1{,}630 & 1.000 & 1{,}630 \\
Census 2011 covariate match    & 1{,}488 & 0.913 & 1{,}488 \\
AOI polygon (after recovery)   & 1{,}549 & 0.950 & 0\textsuperscript{a} \\
District flood extent          & 0       & 0.000 & 0 \\
District news counts           & 0       & 0.000 & 0 \\
\midrule
Final analyzable               & ---     & ---   & \textbf{0} \\
\bottomrule
\end{tabular}
\\[2pt]
\footnotesize\textsuperscript{a}Resolved against the previous registry
fingerprint; invalidated by design when the registry was regenerated.
\end{table}

%=======================================================================
\section{Discussion}
%=======================================================================
The findings above are the kind a pipeline surfaces only if it is allowed to
fail loudly. A system that reused the state corpus would have produced 127{,}252
plausible article rows with no district validity; one that reused the
pre-regeneration satellite cache would have produced 1{,}549 areas measured
against a registry that no longer exists; one without an estimability gate would
have fitted an NB2 model to whatever survived. Each of those is a defensible
engineering shortcut and each would have yielded a publishable-looking
coefficient. Reporting $N=0$ is the cost of refusing all three, and we regard
the refusal, not the count, as the transferable result.

Three limitations bound any future estimate. Census~2011 predates most events,
so urbanization is a fixed pre-period covariate. Boundary references are
mixed-vintage and uncertified against event-date administrative geography. And
GDELT-indexed, district-explicit coverage is a subset of all reporting whose
language and indexing behaviour may itself correlate with urbanization---a
threat the design records rather than solves. Any eventual result would license
statements about coverage at similar \emph{observed flood extent}, not about
intent, causal discrimination or equal total damage.

%=======================================================================
\section{Conclusion}
%=======================================================================
We described CVND, a district-level pipeline joining satellite flood extent,
urbanization and news visibility across 190 Indian flood events, built on the
principle that a measurement must carry the identity of the specification that
produced it. The audit yields four reusable results: boundary recovery is
feasible at 98.7\% but forces a non-additive mixed-vintage geometry; Sentinel-2
L2A cannot support an optical-primary design before 2019, while L1C recovers
about a third of those cases at lower per-scene quality; state-level news
corpora resolve to districts for well under 1\% of in-window articles; and
pre-specified estimability gates will, and should, return \emph{not estimable}
when the inputs do not exist. Completing the district satellite measurement and
district-explicit news collection under the current registry fingerprint is the
immediate next step, after which the pre-registered NB2 specification can be
estimated without modification.

%=======================================================================
\begin{thebibliography}{00}

\bibitem{eisensee2007} T.~Eisensee and D.~Str\"omberg, ``News droughts, news
floods, and U.S. disaster relief,'' \emph{Quarterly Journal of Economics},
vol.~122, no.~2, pp.~693--728, 2007.

\bibitem{emdat} Centre for Research on the Epidemiology of Disasters (CRED),
UCLouvain, ``EM-DAT: The International Disaster Database,'' Brussels, Belgium.

\bibitem{pekel2016} J.-F. Pekel, A.~Cottam, N.~Gorelick, and A.~S. Belward,
``High-resolution mapping of global surface water and its long-term changes,''
\emph{Nature}, vol.~540, pp.~418--422, 2016.

\bibitem{gorelick2017} N.~Gorelick \emph{et al.}, ``Google Earth Engine:
Planetary-scale geospatial analysis for everyone,'' \emph{Remote Sensing of
Environment}, vol.~202, pp.~18--27, 2017.

\bibitem{torres2012} R.~Torres \emph{et al.}, ``GMES Sentinel-1 mission,''
\emph{Remote Sensing of Environment}, vol.~120, pp.~9--24, 2012.

\bibitem{otsu1979} N.~Otsu, ``A threshold selection method from gray-level
histograms,'' \emph{IEEE Trans. Systems, Man, and Cybernetics}, vol.~9, no.~1,
pp.~62--66, 1979.

\bibitem{drusch2012} M.~Drusch \emph{et al.}, ``Sentinel-2: ESA's optical
high-resolution mission for GMES operational services,'' \emph{Remote Sensing of
Environment}, vol.~120, pp.~25--36, 2012.

\bibitem{mcfeeters1996} S.~K. McFeeters, ``The use of the normalized difference
water index (NDWI) in the delineation of open water features,''
\emph{International Journal of Remote Sensing}, vol.~17, no.~7, pp.~1425--1432,
1996.

\bibitem{hilbe2011} J.~M. Hilbe, \emph{Negative Binomial Regression}, 2nd ed.
Cambridge, U.K.: Cambridge University Press, 2011.

\bibitem{census2011} Office of the Registrar General and Census Commissioner,
India, ``Census of India 2011: District-level urban and rural population,'' New
Delhi, India, 2011.

\bibitem{cameron2015} A.~C. Cameron and D.~L. Miller, ``A practitioner's guide
to cluster-robust inference,'' \emph{Journal of Human Resources}, vol.~50,
no.~2, pp.~317--372, 2015.

\bibitem{runfola2020} D.~Runfola \emph{et al.}, ``geoBoundaries: A global
database of political administrative boundaries,'' \emph{PLoS ONE}, vol.~15,
no.~4, e0231866, 2020.

\bibitem{leetaru2013} K.~Leetaru and P.~A. Schrodt, ``GDELT: Global data on
events, location and tone, 1979--2012,'' in \emph{Proc. ISA Annual Convention},
2013.

\bibitem{tellman2021} B.~Tellman \emph{et al.}, ``Satellite imaging reveals
increased proportion of population exposed to floods,'' \emph{Nature}, vol.~596,
pp.~80--86, 2021.

\end{thebibliography}

\end{document}
